% ============================================================================
% Fichier  : partie5/P05_C14_cas_integrateur.tex
% Titre    : Cas Intégrateur : la Mission e-RECETTES
% Auteur   : MINKA MI NGUIDJOÏ Thierry Emmanuel
% Version  : 0.1 -- novembre 2026
% ============================================================================
% J08, 2 h en salle, par équipes de 4 à 5 ; les 2 heures suivantes sont
% offertes. Reprend les faits e-RECETTES de C10, C11, C13. Faits ajoutés :
% profils e-RECETTES : FOMEKONG (ENC, ANN), TAGNE (ENC, ANN, VAL),
% MBIDA (VER), numerika_adm (ADM, prestataire) ; recette totale 2025 :
% 2 184 600 000 FCFA ; versements au Trésor : 2 171 950 000 FCFA (écart
% 12 650 000 FCFA, non expliqué à la date de la mission).
% Corrigé détaillé : livret (corriges/Corr_C14.tex).
% ============================================================================

\chapter{Cas Intégrateur : la Mission e-RECETTES}
\label{chap:cas-integrateur}
\chaptermark{Cas intégrateur}

\epigraph{%
    « Une mission se juge à sa conclusion. »%
}{Règle du cours}

\niveauChapitre{Expert}{L'ensemble du cours.}

\begin{boxobjectifs}
\begin{itemize}
    \item Conduire, en équipe et en temps limité, une mission complète
          d'évaluation du contrôle interne d'un système d'information
          public.
    \item Mobiliser le mandat, la matrice des risques et des contrôles, le
          questionnaire, les tests et le classement des déficiences.
    \item Rédiger au moins deux constats en cinq éléments.
    \item Formuler et défendre oralement une conclusion sur l'efficacité du
          contrôle interne.
\end{itemize}
\end{boxobjectifs}

\bigskip

% ============================================================================
\section*{Mode d'emploi}
% ============================================================================

Le cas se traite en salle, le dernier jour, en deux heures, par équipes de
quatre ou cinq. Chaque équipe est une équipe de vérification du Contrôle
Supérieur de l'État. Elle reçoit le dossier de la section~\ref{sec:c14-dossier},
suit le déroulé de la section~\ref{sec:c14-deroule}, et restitue ses
résultats oralement en dix minutes. Le formateur dispose du corrigé dans
le livret de correction.

\separateur

% ============================================================================
\section{Le mandat}
\label{sec:c14-mandat}
% ============================================================================

\begin{tableaucours}{L{3.0cm}L{9.2cm}}{%
    \tableauHeader{Rubrique} &
    \tableauHeader{Ordre de mission}
}
Fondement & Décret n°~2013/287, art.~2, al.~3 et~4 \\
Objectif & Le contrôle interne d'e-RECETTES garantit-il l'exhaustivité et
la régularité des recettes de la recette des finances de Bafoussam~? \\
Périmètre & e-RECETTES~; le contrat avec NUMÉRIKA~SARL~; les versements
au Trésor \\
Période & Exercice 2025 \\
Livrables & Matrice des risques et des contrôles, cotations, tableau des
déficiences, deux constats, conclusion \\
\end{tableaucours}
\vspace{-0.8em}
\begin{center}{\small\itshape Tableau~14.1. Ordre de mission du cas
intégrateur}\end{center}

\separateur

% ============================================================================
\section{Le dossier}
\label{sec:c14-dossier}
% ============================================================================

\subsection*{Pièce 1. Fiche de l'entité}

La recette des finances de Bafoussam encaisse les recettes de l'État de
son ressort. Depuis le 2~janvier~2025, les encaissements passent par
e-RECETTES, livré par NUMÉRIKA~SARL au titre d'un marché de
640\,000\,000~FCFA. 12\,400 contribuables y sont enregistrés. Trois agents
utilisent l'application.

\subsection*{Pièce 2. Profils d'e-RECETTES}

\begin{tableaucours}{L{3.0cm}L{4.4cm}L{4.8cm}}{%
    \tableauHeader{Compte} &
    \tableauHeader{Fonction} &
    \tableauHeader{Profils}
}
\texttt{fomekong.a} & Caissier & ENC (encaisser), ANN (annuler) \\
\texttt{tagne.m} & Chef de la recette & ENC, ANN, VAL (valider la
journée) \\
\texttt{mbida.s} & Comptable & VER (verser au Trésor) \\
\texttt{numerika\_adm} & Prestataire & ADM (administration) \\
\end{tableaucours}

\subsection*{Pièce 3. Réponses au questionnaire de contrôle interne}

\begin{tableaucours}{L{6.0cm}L{6.2cm}}{%
    \tableauHeader{Question} &
    \tableauHeader{Réponse déclarée}
}
Les quittances sont-elles numérotées automatiquement et sans trou~? &
«~Oui, l'application numérote.~» \\
Une annulation exige-t-elle une validation~? & «~Le chef valide la
journée.~» \\
Les encaissements sont-ils rapprochés des versements~? & «~Chaque fin de
mois.~» \\
Les changements de l'application sont-ils approuvés~? & «~NUMÉRIKA nous
prévient.~» \\
Les sauvegardes sont-elles testées~? & «~NUMÉRIKA s'en occupe.~» \\
\end{tableaucours}

\subsection*{Pièce 4. Résultats d'extractions (faites devant l'équipe)}

\begin{itemize}
    \item Quittances 2025~: Q-000001 à Q-018640~; 37 numéros absents.
    \item Annulations~: 112, dont 61 par \texttt{fomekong.a}~; aucune
          validation d'un second agent n'est enregistrée.
    \item Recettes encaissées en 2025~: 2\,184\,600\,000~FCFA~; versements
          au Trésor~: 2\,171\,950\,000~FCFA.
    \item Aucune trace d'un rapprochement mensuel.
\end{itemize}

\subsection*{Pièce 5. Journaux techniques}

\begin{itemize}
    \item 14 mises en production en 2025, dont 11 sans demande écrite.
    \item 9~septembre~2025~: modification directe du taux de pénalité par
          \texttt{numerika\_adm}.
    \item Sauvegardes~: 23 échecs, dont 6 consécutifs~; aucune
          restauration testée.
\end{itemize}

\subsection*{Pièce 6. Contrat et réception}

\begin{itemize}
    \item Ni propriété des données, ni réversibilité, ni droit d'audit~;
          hébergement hors du pays.
    \item Procès-verbal de réception du 15~janvier~2025~: 47 anomalies
          ouvertes, dont 9 bloquantes.
    \item Aucun audit de sécurité de l'ANTIC.
\end{itemize}

\subsection*{Pièce 7. Réponses des répondants (instrument OUI / NON / NEK)}

Cinq questions par domaine, posées à quatre fonctions (chef de la recette,
caissier, comptable, correspondant de NUMÉRIKA), soit vingt réponses par
domaine.

\begin{tableaucours}{L{5.0cm}ccc}{%
    \tableauHeader{Domaine} &
    \tableauHeader{\OUI{}} & \tableauHeader{\NON{}} & \tableauHeader{\NEK{}}
}
R1 Encaissement, quittances, annulations & 9 & 4 & 7 \\
R2 Accès et changements & 6 & 7 & 7 \\
R3 Sauvegardes et continuité & 3 & 9 & 8 \\
\end{tableaucours}

\separateur

% ============================================================================
\section{Le déroulé}
\label{sec:c14-deroule}
% ============================================================================

\begin{tableaucours}{L{2.0cm}L{5.4cm}L{4.8cm}}{%
    \tableauHeader{Temps} &
    \tableauHeader{Travail de l'équipe} &
    \tableauHeader{Chapitres mobilisés}
}
0 -- 10 min & Lire le mandat~; répartir les rôles (chef de mission,
rédacteur, deux testeurs) & \ref{chap:institutions},
\ref{chap:preparer} \\
10 -- 30 min & Construire la matrice des risques et des contrôles
(au moins 10 lignes) & \ref{chap:eval-sci}, \ref{chap:applications} \\
30 -- 50 min & Confronter le QCI aux pièces~; coter chaque contrôle~;
calculer $\mu$, $U$, $H$ (pièce~7) et situer les domaines dans le quadrant
& \ref{chap:preuve}, \ref{chap:constater} \\
50 -- 65 min & Regrouper et classer les déficiences & \ref{chap:eval-sci} \\
65 -- 85 min & Rédiger deux observations et la conclusion &
\ref{chap:constater} \\
85 -- 105 min & Restitution~: 5 à 7 minutes par équipe, questions du jury
& tout le cours \\
105 -- 120 min & Ouverture sur la recherche~: quel sujet de mémoire ce
cas suggère-t-il~? & \ref{chap:recherche} \\
\end{tableaucours}
\vspace{-0.8em}
\begin{center}{\small\itshape Tableau~14.2. Déroulé du cas en deux
heures}\end{center}

\separateur

% ============================================================================
\section{La grille de restitution}
\label{sec:c14-grille}
% ============================================================================

\begin{tableaucours}{L{7.8cm}c}{%
    \tableauHeader{Critère} &
    \tableauHeader{Points}
}
Le mandat est cité et l'objectif reformulé correctement & 2 \\
La matrice part des risques et couvre accès, changements, exploitation,
application, versements, prestataire & 4 \\
Les cotes distinguent les déclarations (QCI) des preuves~; les trois
domaines sont situés dans le quadrant & 3 \\
Le classement identifie au moins une déficience majeure, et la justifie &
4 \\
Les deux observations ont leurs sept éléments et parlent de comptes, pas
de personnes & 4 \\
La conclusion porte sur un objectif, une période et un périmètre & 3 \\
\textbf{Total} & \textbf{20} \\
\end{tableaucours}
\vspace{-0.8em}
\begin{center}{\small\itshape Tableau~14.3. Grille d'évaluation de la
restitution}\end{center}

\separateur

% ============================================================================
\section{Les erreurs fréquentes}
\label{sec:c14-erreurs}
% ============================================================================

\begin{tableaucours}{L{5.8cm}L{6.4cm}}{%
    \tableauHeader{Erreur} &
    \tableauHeader{Ce qu'il fallait faire}
}
Conclure « conforme à 45~\% » sur R1 & $U = 0{,}35 \geq 0{,}20$~: le
domaine est opaque~; aucune conclusion de conformité (règle~1) \\
Coter \OUI{} la numérotation parce que l'application «~numérote~» &
Partir des 37 numéros absents~: \NON{} \\
Coter \NEK{} le rapprochement mensuel & Aucune trace~: l'absence d'une
trace exigée est un \NON{} \\
Écrire «~M.~FOMEKONG détourne des recettes~» & Constater la
concentration des annulations sur un compte, sans validation \\
Conclure sur la sécurité d'ensemble & Rester dans l'objectif~; signaler
l'absence d'audit de l'ANTIC \\
Additionner les \NON{} comme autant de déficiences & Regrouper par cause
(accès, changements, projet) \\
Oublier l'écart de 12\,650\,000~FCFA & C'est l'effet chiffré le plus
direct sur l'objectif d'exhaustivité \\
\end{tableaucours}
\vspace{-0.8em}
\begin{center}{\small\itshape Tableau~14.4. Grille des erreurs
fréquentes}\end{center}

\begin{boxmetier}{Vérificateur}
Dans un vrai dossier, l'écart entre recettes encaissées et versements
appellerait une note de faits distincte du rapport
(Chapitre~\ref{chap:constater})~: c'est une suite possible devant le CDBF.
\end{boxmetier}

\begin{boxmetier}{Juge des comptes}
L'écart entre encaissements et versements concerne directement le compte
du comptable~: il sera examiné, par extension, à travers les comptes
produits par e-RECETTES.
\end{boxmetier}

\separateur

\begin{boxresume}
\begin{itemize}
    \item \textbf{Une mission complète} tient en six gestes~: mandat,
          matrice, cotation, classement, constats, conclusion.
    \item \textbf{Le QCI se confronte aux pièces}~; une déclaration ne vaut
          pas preuve.
    \item \textbf{Les déficiences se regroupent par cause.}
    \item \textbf{La conclusion se défend}~: objectif, période, périmètre,
          travaux.
\end{itemize}
\end{boxresume}

\begin{boxexo}[title={\ding{45}\quad Exercice~14.1. La matrice \niveaumoyen}]
Construisez la matrice des risques et des contrôles d'e-RECETTES (au moins
dix lignes), à partir des pièces du dossier.
\end{boxexo}

\begin{boxexo}[title={\ding{45}\quad Exercice~14.2. Les déficiences \niveauexpert}]
Regroupez les contrôles cotés \NON{} en déficiences, classez-les avec la
grille du Chapitre~\ref{chap:eval-sci}, et justifiez chaque déficience
majeure.
\end{boxexo}

\begin{boxexo}[title={\ding{45}\quad Exercice~14.3. La conclusion \niveauexpert}]
Rédigez la conclusion de la mission et le premier paragraphe de la
synthèse du rapport.
\end{boxexo}

\begin{boxinfo}
\textbf{Transition.} La mission est complète. Le
Chapitre~\ref{chap:recherche} remonte aux fondements de la méthode, le
programme de recherche de l'auditologie, et ouvre des pistes de mémoire
pour les auditeurs qui voudront contribuer à le tester.
\end{boxinfo}
